\documentclass[a4paper,12pt]{article}
\usepackage[T2A]{fontenc}
\usepackage[utf8]{inputenc}
\usepackage[russian]{babel}
\usepackage{amsmath,amssymb,mathtools}
\usepackage{PTSans}
\renewcommand{\familydefault}{\sfdefault}
\usepackage{icomma}
\usepackage[a4paper,margin=18mm,headheight=25pt]{geometry}
\usepackage{microtype}
\usepackage{tikz}
\usetikzlibrary{calc,arrows.meta,positioning,shapes.geometric}
\usepackage{tabularx,booktabs,longtable}
\usepackage{enumitem}
\usepackage[hidelinks]{hyperref}
\usepackage{parskip}
\usepackage{fancyhdr}
\usepackage{setspace}
\usepackage{xcolor}

\definecolor{accent}{HTML}{176B6B}
\definecolor{darkaccent}{HTML}{0E4747}
\definecolor{textgray}{HTML}{2F3437}
\definecolor{softgray}{HTML}{6A7377}
\definecolor{rulegray}{HTML}{CFD8D8}

\newcommand{\LessonDate}{29.09.26}
\onehalfspacing
\setlength{\parindent}{0pt}
\setlist[itemize]{leftmargin=6mm,itemsep=1pt,topsep=2pt}
\setlist[enumerate]{leftmargin=8mm,itemsep=4pt,topsep=4pt}

\pagestyle{fancy}
\fancyhf{}
\fancyhead[L]{\small\color{softgray}Уравнения с модулем}
\fancyhead[R]{\small\color{softgray}\LessonDate}
\fancyfoot[C]{\small\color{softgray}\thepage}
\renewcommand{\headrulewidth}{0.3pt}
\renewcommand{\headrule}{\hbox to\headwidth{\color{rulegray}\leaders\hrule height \headrulewidth\hfill}}

\newcommand{\topic}[1]{%
  \vspace{0.5em}\noindent{\Large\bfseries\color{darkaccent}#1}\par
  \vspace{0.25em}{\color{rulegray}\hrule}\vspace{0.55em}}
\newcommand{\subtopic}[1]{\vspace{0.3em}\noindent{\large\bfseries\color{accent}#1}\par\vspace{0.2em}}
\newcommand{\key}[1]{\textbf{\color{darkaccent}#1}}
\newcommand{\R}{\mathbb{R}}

\begin{document}

% ---------- TITLE PAGE ----------
\thispagestyle{empty}
\begin{center}
\vspace*{20mm}
{\Large\color{softgray}10 класс}\par
\vspace{13mm}
{\fontsize{32}{37}\selectfont\bfseries\color{darkaccent}Чек-лист}\par
\vspace{5mm}
{\LARGE\bfseries Уравнения с модулем}\par
\vspace{3mm}
{\large знаки, интервалы и равносильные переходы}\par
\vspace{14mm}
{\large Анна Трапезникова}\par
\vspace{3mm}
{\normalsize \LessonDate}\par
\vfill
{\small\color{softgray}Лёвин Артём Александрович эксклюзивно для Анны Трапезниковой}\par
\vspace{7mm}
\begin{tikzpicture}[x=1cm,y=1cm]
  \draw[accent!60,line width=1.2pt]
    (-7,0) .. controls (-4.5,1.0) and (-1.7,-0.9) .. (0,0)
           .. controls (2.3,1.15) and (4.7,-0.8) .. (7,0.25);
\end{tikzpicture}
\end{center}
\clearpage

% ---------- MAIN CHECKLIST ----------
\topic{1. Главная идея модуля}

Для любого выражения $F(x)$:
\[
|F(x)|=
\begin{cases}
F(x), & F(x)\ge 0,\\
-F(x), & F(x)<0.
\end{cases}
\]

\key{Ключевой вопрос перед раскрытием модуля:} какой знак имеет подмодульное выражение на рассматриваемых значениях $x$?

\subtopic{Простейшая модель}
\[
|x+2|=1.
\]
Можно рассмотреть два случая:
\[
\left\{
\begin{aligned}
x+2&\ge0,\\
x+2&=1
\end{aligned}
\right.
\qquad\text{или}\qquad
\left\{
\begin{aligned}
x+2&<0,\\
-(x+2)&=1.
\end{aligned}
\right.
\]
Полученные корни обязательно проверяются на принадлежность своему случаю.

\subtopic{Быстрый частный случай}
Если справа стоит постоянное число $a$, то:
\[
|F(x)|=a.
\]
При $a<0$ решений нет; при $a=0$ решаем $F(x)=0$; при $a>0$ можно сразу перейти к
\[
F(x)=a \quad\text{или}\quad F(x)=-a.
\]
Это тот же смысл раскрытия модуля, записанный короче.

\topic{2. Один модуль слева: равносильный переход}

Если уравнение приведено к виду
\[
|F(x)|=G(x),
\]
то правая часть обязана быть неотрицательной. Поэтому
\[
|F(x)|=G(x)
\iff
\left[
\begin{array}{l}
\left\{
\begin{aligned}
G(x)&\ge0,\\
F(x)&=G(x)
\end{aligned}
\right.\\[2mm]
\left\{
\begin{aligned}
G(x)&\ge0,\\
F(x)&=-G(x)
\end{aligned}
\right.
\end{array}
\right.
\]

\key{Практический приём:} часто условие $G(x)\ge0$ удобнее использовать как фильтр для найденных кандидатов, чем отдельно решать сложное неравенство.

\subtopic{Пример}
\[
|2x-1|=x+5.
\]
Сначала учитываем $x+5\ge0$. Затем решаем
\[
2x-1=x+5 \quad\text{или}\quad 2x-1=-(x+5).
\]
Получаем $x=6$ и $x=-\frac43$; оба удовлетворяют условию $x+5\ge0$.

\key{ОДЗ остаётся обязательной.} Если в $F(x)$ или $G(x)$ есть знаменатели, корни или другие ограничения, их проверяют независимо от условия $G(x)\ge0$.

\topic{3. Несколько модулей: метод интервалов}

Когда в уравнении несколько модулей, обычно надёжнее работать через числовую ось.

\subtopic{Алгоритм}
\begin{enumerate}
  \item Выпишите все подмодульные выражения.
  \item Найдите точки, где каждое из них обращается в ноль.
  \item Отметьте эти точки на числовой оси: они разбивают её на интервалы.
  \item На каждом интервале определите знак каждого подмодульного выражения подстановкой любой удобной точки.
  \item Раскройте все модули согласно найденным знакам.
  \item Решите полученное обычное уравнение.
  \item Оставьте только корни, принадлежащие рассматриваемому интервалу.
  \item Объедините результаты всех случаев.
\end{enumerate}

\subtopic{Пример}
Решим
\[
|x-2|+|x+1|=7.
\]
Критические точки: $x=-1$ и $x=2$.

\begin{center}
\begin{tikzpicture}[x=2cm,y=1cm,>=Stealth]
  \draw[->,line width=0.9pt] (-2.4,0) -- (2.6,0) node[right] {$x$};
  \foreach \x/\lab in {-1/-1,1/2}{
    \draw[accent,line width=1pt] (\x,-0.11)--(\x,0.11);
    \node[below=3pt] at (\x,0) {$\lab$};
  }
  \node[above=5pt] at (-1.7,0) {$(-,-)$};
  \node[above=5pt] at (0,0) {$(-,+)$};
  \node[above=5pt] at (1.8,0) {$(+,+)$};
  \node[font=\small,softgray,below=18pt] at (0,0) {знаки $(x-2,\;x+1)$};
\end{tikzpicture}
\end{center}

На трёх интервалах получаем три обычных уравнения. После проверки принадлежности интервалам остаются корни $x=-3$ и $x=4$.

\subtopic{Особенно внимательно со знаком перед модулем}
Если есть выражение $-|F(x)|$ и на выбранном интервале $F(x)<0$, то
\[
-|F(x)|=-(-F(x))=F(x).
\]
Два минуса должны быть обработаны последовательно.

\topic{4. Как выбрать способ}

\begin{itemize}
  \item \key{Один изолированный модуль слева:} обычно выгоден переход $|F|=G$ с условием $G\ge0$.
  \item \key{Несколько модулей в сумме или разности:} обычно удобнее числовая ось и разбиение на интервалы.
  \item \key{Вложенный модуль:} часто удобно последовательно уединять внешний, затем внутренний модуль.
  \item \key{Сложное подмодульное выражение:} прежде чем строить много случаев, проверьте, можно ли уединить модуль и применить равносильный переход.
\end{itemize}

\topic{5. Типичные ошибки}

\begin{itemize}
  \item раскрыть $|F(x)|$ без выяснения знака $F(x)$;
  \item забыть условие $G(x)\ge0$ в уравнении $|F(x)|=G(x)$;
  \item потерять ОДЗ дроби, даже если остальные условия выполнены;
  \item решить уравнение для одного интервала и не проверить, принадлежит ли найденный корень этому интервалу;
  \item неправильно обработать внешний минус перед модулем;
  \item считать, что каждый новый интервал обязан дать корень: иногда получается неверное равенство и решений в этом случае нет;
  \item забыть, что на границе интервалов подмодульное выражение равно нулю и такая точка может быть решением;
  \item выбирать громоздкий способ, когда структура уравнения допускает более короткий метод.
\end{itemize}

% ---------- FLOWCHART PAGE ----------
\clearpage
\thispagestyle{plain}
\begin{center}
{\Large\bfseries\color{darkaccent}Алгоритм: выбор метода для уравнения с модулем}\par
\vspace{10mm}
\begin{tikzpicture}[
  node distance=10mm and 14mm,
  every node/.style={font=\small,align=center},
  start/.style={ellipse,draw=accent,line width=1pt,minimum width=37mm,minimum height=10mm},
  action/.style={rectangle,rounded corners=2mm,draw=accent,line width=1pt,minimum width=49mm,minimum height=12mm,text width=45mm},
  decision/.style={diamond,aspect=2.0,draw=accent,line width=1pt,text width=37mm,inner sep=2pt},
  arr/.style={-{Stealth[length=2.5mm]},line width=0.9pt,color=darkaccent}
]
  \node[start] (s) {Дано уравнение с модулем};
  \node[action,below=of s] (prep) {Упростите выражения и проверьте ОДЗ};
  \node[decision,below=of prep] (one) {Можно изолировать один модуль слева?};
  \node[action,below left=14mm and 15mm of one] (eqv) {Приведите к $|F(x)|=G(x)$; учтите $G(x)\ge0$; решите $F=G$ и $F=-G$};
  \node[action,below right=14mm and 15mm of one] (axis) {Найдите нули всех подмодульных выражений и разбейте ось на интервалы};
  \node[action,below=20mm of axis] (signs) {На каждом интервале определите знаки и раскройте модули};
  \node[action,below=20mm of eqv] (check1) {Проверьте кандидаты по ОДЗ и условию $G(x)\ge0$};
  \node[action,below=of signs] (check2) {Проверьте принадлежность каждого корня своему интервалу};
  \coordinate (merge) at ($(check1.south)!0.5!(check2.south)+(0,-18mm)$);
  \node[start,below=7mm of merge] (end) {Объедините допустимые решения};

  \draw[arr] (s)--(prep);
  \draw[arr] (prep)--(one);
  \draw[arr] (one)-- node[above left]{да} (eqv);
  \draw[arr] (one)-- node[above right]{нет} (axis);
  \draw[arr] (eqv)--(check1);
  \draw[arr] (axis)--(signs);
  \draw[arr] (signs)--(check2);
  \draw[arr] (check1.south) |- (merge);
  \draw[arr] (check2.south) |- (merge);
  \draw[arr] (merge) -- (end.north);
\end{tikzpicture}
\end{center}
\clearpage

% ---------- TRAINING ----------
\topic{Тренировочный блок — Путь А}

Решите уравнения. В заданиях 3, 4 и 9 сначала приведите уравнение к виду, удобному для выбранного метода. В задачах с несколькими модулями явно отметьте критические точки на оси.

\begin{enumerate}
  \item $|3x-5|=7$.
  \item $|x^2-4|=5$.
  \item $|2x-1|=x+5$.
  \item $\displaystyle \left|\frac{2x-3}{x+1}\right|=\frac{x+5}{x+1}$.
  \item $|x-2|+|x+1|=7$.
  \item $|2x+3|-|x-4|=2$.
  \item $|x-1|+|x+2|=3$.
  \item $|x-4|-|x+2|=6$.
  \item $|x|=2+|x-3|$.
  \item $\bigl||x-2|-3\bigr|=1$.
\end{enumerate}

\vfill
{\small\color{softgray}\textbf{Самопроверка перед ответом:} учтены ли ОДЗ; проверены ли знаки; принадлежит ли корень своему интервалу; не потеряна ли граница; выполнено ли условие неотрицательности правой части.}

\clearpage

% ---------- ANSWERS ----------
\topic{Ответы}

\begin{center}
\renewcommand{\arraystretch}{1.35}
\begin{tabularx}{0.92\linewidth}{@{}>{\bfseries}cX@{}}
\toprule
№ & \textbf{Ответ} \\
\midrule
1 & $x=-\dfrac23,\;4$ \\
2 & $x=-3,\;3$ \\
3 & $x=-\dfrac43,\;6$ \\
4 & $x=-\dfrac23,\;8$ \\
5 & $x=-3,\;4$ \\
6 & $x=-9,\;1$ \\
7 & $x\in[-2;1]$ \\
8 & $x\in(-\infty;-2]$ \\
9 & $x=\dfrac52$ \\
10 & $x=-2,\;0,\;4,\;6$ \\
\bottomrule
\end{tabularx}
\end{center}

\vspace{8mm}
\subtopic{Контрольные идеи}
\begin{itemize}
  \item В №4 обязательно исключается $x=-1$ по ОДЗ.
  \item В №7 множество решений является целым отрезком: на $[-2;1]$ сумма расстояний до точек $-2$ и $1$ постоянна и равна $3$.
  \item В №8 решением является луч $(-\infty;-2]$; это нормальный результат для уравнения с модулями.
  \item В №10 внешний модуль сначала сводит задачу к двум значениям внутреннего выражения, после чего раскрывается внутренний модуль.
\end{itemize}

\vfill
\begin{center}
{\small\color{softgray}Лёвин Артём Александрович эксклюзивно для Анны Трапезниковой}
\end{center}

\end{document}